\documentclass[../../main.tex]{subfiles}
\graphicspath{{\subfix{../../image/}}}

\begin{document}

\section*{第二章小结}\label{section:ch02-第二章小结}

\addcontentsline{toc}{section}{第二章小结}

本章中我们讨论了（统计）独立性的概念, 它是概率论中最重要的概念之一. 独立性是概率论特有的概念, 它的引进大大推动了概率论的发展, 前期概率论中最重要的一些结果大都是在独立性的假定下获得的, 只有到了近代才开始研究一些不独立但常在另一种较弱独立性假定下的概率模型. 读者能在本书的后面部分充分体验到独立性的重要性.

我们定义了事件的独立性与试验的独立性, 读者应把它们进行对比. 此外, 两个事件独立性的定义与多个事件独立性的定义形式上也有区别, 应理解为什么会发生这种情况.

条件概率是概率论中另一重要概念, 它与独立性有密切联系. 在不具有独立性的场合, 它将扮演重要角色. 条件概率也是某种概率.

本章还导出几个基本公式：乘法公式及其推广形式, 全概率公式, Bayes 公式, 它们以后经常被用到.

Bernoulli 试验是概率论中最重要的模型之一. 正是通过对这个模型的不断深入研究, 逐渐提出了概率论特有的课题, 创造出相应的工具与方法. 它对后来的发展有着不可估量的影响. Bernoulli 试验概型在应用上也很重要. 在第五章, 我们将继续讨论这个概型, 证明有关的极限定理并最后解决一些本章遗留的计算问题.

随机游动是概率论中最早研究的一个动态模型, 可追溯到 17 世纪 Huygens 对公平博弈中赌徒输光问题的研究, 它开了随机过程理论研究的先河. 当代, 随机游动理论的研究仍在深入, 应用既广且深.

二项分布与 Poisson 分布是概率论中最重要的三个分布中的两个, 另一个分布——正态分布, 将在下章出现. 这两个分布在理论研究和实际应用中都很重要. 二项分布广泛应用于抽样检查等场合, 而 Poisson 分布则大量出现于社会生活和物理现象中. Poisson 过程的结构分析是随机过程理论的最基本成果之一.

从某些假定出发, 利用全概率公式导出某个方程, 最后通过解方程求出答案, 在随机游动和 Poisson 过程研究中使用的这种方法颇为有力, 它是随机过程理论研究中的分析方法的特例.

到目前为止, 本课程的展开基本上还是与概率论的历史发展相平行的, 而且, 我们已经为概率论另一重要概念——随机变量的引入作了不少准备.

\end{document}